\section{Closed tests and the finite-observation reduction}\label{sec:category}

We first separate the topological argument from the evolution estimates.
The facts in this section are independent of the Kerr parameters. We
include their short proofs to make precise why a finite tidal
observation suffices, and why no regular endpoint tangent is part of
the extension hypothesis.

\begin{lemma}[The data topology]\label{cat:baire}
The space \(\mathcal D\) is completely metrizable and
\(\mathcal U_\epsilon\) is Baire. Every neighborhood of a datum contains
a ball \(p_m(d'-d)<\eta\) relative to \(\mathcal D\), for some finite
\(m\) and \(\eta>0\).
\end{lemma}
\begin{proof}
A sequence Cauchy in every \(p_m\) has compatible smooth limits on
compact coordinate sets. Passing to the pointwise limit in its weighted
uniform Cauchy bounds proves convergence to that limit in every
\(p_m\), including the finiteness of each limiting seminorm. Thus
\(E\) is a Fr\'echet space. For example its topology is defined by
the complete metric
\(\sum_{m\geq0}2^{-m-1}\min\{1,p_m(d-d')\}\).

Any positive metric in the affine space is uniformly comparable to
\(h_*\): this follows from the decaying difference on the ends and
compactness on their complement. Hence it is complete, and positivity
is open already in the \(p_0\) topology. The constraint equations are
closed under local smooth convergence within this positive open set.
An open subset of a completely metrizable space is completely
metrizable, as is a relatively closed subset thereof. Consequently
\(\mathcal D\), and then its open subset \(\mathcal U_\epsilon\),
are completely metrizable and Baire. Finally the seminorms increase
with \(m\), so one sufficiently small ball in the largest of finitely
many orders lies inside any basic neighborhood.
\end{proof}

\subsection{Seeds and curvature measurements}

A seed \(w\in T\Sigma\) includes its base point. For datum
\(d=(h,K)\), its initial future unit velocity is
\[
 U_d(w)=\sqrt{1+h(w,w)}\,n+w,
\]
where \(n\) is the future unit normal. Let
\(\gamma_{d,w}:[0,T_d(w))\to M_d\) be the maximal future geodesic
with this initial velocity, parametrized by proper time. Along it
parallel transport the positive inner product \(1\oplus h\) on
\(\mathbb R n\oplus T\Sigma\), and call the result \(k_{d,w}\).
Set
\begin{equation}\label{cat:Q}
 Q_d(w,\tau)=|\operatorname{Riem}(g_d)|_{k_{d,w}(\tau)}.
\end{equation}
The curvature tensor is covariant and the norm is its positive
Hilbert--Schmidt norm, independently of any initial orthonormal frame.

Choose a countable relatively compact open basis \(\mathscr O\) of
\(T\Sigma\). For \(O\in\mathscr O\) and an integer \(N\geq1\), define
\begin{equation}\label{cat:test}
 A(O,N)=\left\{d\in\mathcal D:
 \begin{array}{ll}
 T_d(w)\leq N &\text{for every }w\in O,\\
 Q_d(w,\tau)\leq N &\text{for every }w\in O,
                    \ 0\leq\tau<T_d(w)
 \end{array}\right\}.
\end{equation}

\begin{lemma}[Closed tests]\label{cat:closed-tests}
Each set \(A(O,N)\) is closed in \(\mathcal D\).
\end{lemma}
\begin{proof}
Every strict failure has a compact witness. If \(T_d(w)>N\), use a
geodesic segment extending beyond \(N\). If a curvature bound fails,
use a segment through a point \(\tau<T_d(w)\) with
\(Q_d(w,\tau)>N\). Smooth local Cauchy stability and geometric
uniqueness place a neighborhood of this fixed compact segment in a
common gauge for all sufficiently close data. For the compact segment
\(K\), its causal shadow \(J^-(K)\cap\Sigma\) is compact; otherwise a limit of
past causal curves to escaping initial points would contradict
Cauchyness. Finite local Cauchy constructions therefore suffice.
Geodesic and parallel-transport ODE stability preserve the strict
witness. This proves openness of the complement, without asserting
continuity of the maximal duration.
\end{proof}

\begin{lemma}[Tidal contractions]\label{cat:tidal}
If \(U=\dot\gamma_{d,w}\), \(g_d(X,X)=1\), and \(g_d(U,X)=0\), then
\begin{equation}\label{cat:tidal-bound}
 |\operatorname{Riem}(U,X,U,X)|
 \leq (1+2h(w,w))^2 Q_d(w,\tau).
\end{equation}
The factor is uniformly bounded for data and seeds near a fixed pair.
\end{lemma}
\begin{proof}
In the transported Lorentz orthonormal frame write
\(U=\alpha e_0+v\), \(X=\beta e_0+z\). Orthogonality and the unit
conditions give \(\alpha^2=1+|v|^2\), \(|z|^2=1+\beta^2\), and
\(\alpha\beta=v\cdot z\); hence \(\beta^2\leq|v|^2\).
Therefore \(|X|_k^2\leq|U|_k^2=1+2h(w,w)\), the latter identity
being constant by parallel transport. Tensor Cauchy--Schwarz proves
the assertion.
\end{proof}

\subsection{From an arbitrary extension to regular exits}

We will use an intrinsic compactness fact. It avoids imposing causal
convexity on the image in an extension.

\begin{lemma}[Fixed-endpoint homotopy]\label{cat:homotopy}
Let \(P\) be a globally hyperbolic open subspacetime of a Hausdorff
Lorentzian manifold. In a continuous family of future causal paths
with fixed endpoints \(p,q\in P\), if the first path lies in \(P\),
then all paths lie in \(P\).
\end{lemma}
\begin{proof}
Parametrize the family by a connected compact interval. The set of
parameters for which the whole path lies in \(P\) is open. Every
such path lies in the intrinsic compact diamond
\(J_P^+(p)\cap J_P^-(q)\), which is closed in the ambient Hausdorff
manifold. Continuity makes the parameter set closed as well. It is
therefore the whole interval.
\end{proof}

\begin{proposition}[Regular-exit covering]\label{cat:regular-exit}
If \(d\in\mathcal U_\epsilon\) admits a \(C^2\) future extension,
there exist \(O\in\mathscr O\) and \(N\geq1\) such that
\(d\in A(O,N)\) and every seed in \(O\) reaches its finite MGHD
endpoint on a regular ambient geodesic arc. Here regular means that
position and unit tangent continue to and through the endpoint.
\end{proposition}
\begin{proof}
Starting from \Cref{intro:extension}, identify the development with its image and
write \(z=c(1)\). The curve is future-inextendible within the
development: an interior endpoint would equal \(z\) by the Hausdorff
property. Extend it maximally to the past inside the development.
Cauchyness then puts a sufficiently late point \(p=c(s_0)\) to
the intrinsic future of \(\Sigma\). In a small temporal chart at
\(z\), the local timelike relation from \(p\) to \(z\) gives a
short timelike geodesic continuing past \(z\). Only this local
relation is needed, not a continuing unit tangent of \(c\).
For a \(C^2\) metric the Christoffel symbols are \(C^1\), so the
geodesic flow has the usual local uniqueness and continuous dependence.

Consider sufficiently close geodesic arcs with starting positions on
the same temporal slice inside the image, and close unit directions.
They remain in a common compact chart and phase-space neighborhood.
Parametrize these arcs by chart time on a common interval
\([t_0,t_1]\), with the reference arc passing through \(z\) at
\(t_z\in(t_0,t_1)\). Each nearby arc must exit the image. To see this,
suppose an arc \((t,x_\gamma(t))\) stayed inside. Choose a fixed smooth
bump \(\beta\), supported in \((t_0,t_1)\), with \(\beta(t_z)=1\).
The paths
\[
 \bigl(t,x_\gamma(t)+\lambda\beta(t)
                    (x(z)-x_\gamma(t_z))\bigr),\qquad 0\le\lambda\le1,
\]
form a fixed-endpoint timelike homotopy: after shrinking the arc family,
the spatial change is uniformly small in \(C^1\), and all paths stay
in the chart with its uniform timelike margin. Under the supposition
that the original arc stayed inside, both endpoints belong to the image.
\Cref{cat:homotopy} would then keep every path inside, although the
path at \(\lambda=1\) passes through \(z\), a contradiction.
Thus every nearby arc has a first exit, with position
and tangent regular there. By geodesic uniqueness this first exit is
its maximal intrinsic endpoint, regardless of possible later reentry.

The reference arc continued backwards inside the development meets
\(\Sigma\). To see that a maximal timelike geodesic cannot end at
an interior position limit, use a precompact temporal chart
\((t,x)\). Causal slopes \(V=(1,\mathrm dx/\mathrm dt)\) are
bounded and the geodesic equation gives
\[
 \frac{\mathrm d}{\mathrm dt}\log\frac{\mathrm dt}{\mathrm d\tau}
 =-\Gamma^0_{\alpha\beta}V^\alpha V^\beta.
\]
Over a finite chart-time tail this bounds the unit tangent and gives
its limit, so the geodesic continues. Cauchyness now gives the claimed
intersection. Flow dependence through that fixed earlier compact
segment yields an open seed family whose arcs have the exits just
constructed.

Shrink the seed family to a relatively compact one. The earlier tube
has uniformly bounded duration, parallel frames and curvature. In the
extension chart, unit tangents stay in a fixed compact phase-space
set. The parallel-transport ODE has bounded coefficients on a bounded
chart-time interval, so its transport matrices and their inverses are
bounded; curvature is continuous and bounded as well. Duration and the norm
\eqref{cat:Q} are consequently uniformly bounded up to first exit.
Choose a basis set \(O\) within the seed family and an integer \(N\)
larger than both bounds.
\end{proof}

Let \(B_\epsilon\) denote the extendible data in
\(\mathcal U_\epsilon\). Define \(\mathcal B(O,N)\) to consist of
the data in \(B_\epsilon\cap A(O,N)\) for which at least one seed
in \(O\) has a regular exit in some extension. Then
\begin{equation}\label{cat:cover}
 B_\epsilon=\bigcup_{O\in\mathscr O,\,N\geq1}\mathcal B(O,N),
 \qquad
 \overline{\mathcal B(O,N)}^{\,\mathcal U_\epsilon}
 \subset A(O,N)\cap\mathcal U_\epsilon.
\end{equation}
The pieces \(\mathcal B(O,N)\) themselves need not be closed.

\subsection{The analytic task}

\begin{proposition}[Finite-observation criterion]\label{cat:destruction}
Suppose a single \(\epsilon>0\) has the following property. For every
\(O\in\mathscr O\), \(N\geq1\), \(d\in\mathcal B(O,N)\),
integer \(m\geq10\) and \(\eta>0\), there exist
\(d'\in\mathcal U_\epsilon\), \(w'\in O\), and
\(0\leq\tau<T_{d'}(w')\) such that
\[
 p_m(d'-d)<\eta,\qquad Q_{d'}(w',\tau)>N.
\]
Then \Cref{thm:main} holds for this \(\epsilon\).
\end{proposition}
\begin{proof}
If the relative closure of \(\mathcal B(O,N)\) had nonempty interior,
that interior would meet the piece and, by \eqref{cat:cover}, be
contained in \(A(O,N)\). At an intersection point, choose a finite
seminorm ball within the interior, increasing the order to at least
ten if necessary. The asserted perturbation is in that ball but fails
the test, a contradiction. Hence each closure is nowhere dense. The
countable intersection of their relatively open dense complements
contains no extendible datum, by \eqref{cat:cover}.
\end{proof}

All choices needed to fix the neighborhood in
\Cref{cat:destruction} will precede the choice of
\(O,N,m,\eta\). The exterior and interior estimates construct the
background comparison region at this common smallness threshold.
Higher constants may then depend on the fixed smooth datum. Only
afterwards do we choose the late finite perturbation that violates
the selected test.
